% ------------------------------------------------------------------------------
% §2 Institutional setting — drafted Phase C from:
%   Direction_clean/README.md; findings_task1b.md (gross payment, post-fix);
%   F1/F2b readout; memory: project_sa_regime_exit_2025, project_aemo_qed_direction_cost
% ------------------------------------------------------------------------------
\section{Institutional Setting}\label{sec:setting}

\subsection{Directions in the National Electricity Market}

The Australian National Electricity Market (NEM) is an energy-only gross pool:
generators submit offer curves --- up to ten price bands with quantities revisable
until dispatch\footnote{Rebidding is regulated: since 2016, offers must not
be false or misleading and late rebids require contemporaneous records
\citep{AEMC2015_goodfaith}. The conduct studied in this paper operates
through availability declarations and band quantities, which are and remain
freely revisable.} --- and a security-constrained dispatch engine clears the market
every five minutes. South Australia is the NEM's high-renewables frontier. Wind
and rooftop solar routinely supply most of regional demand, and on sunny, windy
days the energy market has no commercial work for the state's gas-fired
synchronous generators.

The energy market clears energy; it does not clear system security. South
Australia's grid requires a minimum complement of synchronous machines online to
provide system strength and inertia, codified as a set of minimum synchronous
generator combinations that the Australian Energy Market Operator (AEMO) must
maintain. When market dispatch leaves the requirement unmet --- typically in
exactly those high-renewables, low-price periods --- AEMO issues a
\emph{direction}: an instruction to a specific generating unit to synchronise,
or to remain synchronised, outside the normal dispatch process. A direction
compels presence, not output. A directed unit must be online; it typically runs
at its minimum stable loading.\footnote{Section~\ref{sec:mechanism} measures
this directly: directed output equals the unit's operating floor almost
exactly, with 55.9\% of resolved direction episodes delivering zero or negative
output in excess of the floor block.}

\subsection{Compensation: a gross payment at a backward-looking reference price}

A directed participant is entitled to compensation under clause 3.15.7 of
the National Electricity Rules (NER), with the cost recovered from market
participants under clause 3.15.8 \citep{AEMC2024_ner}. The reference price --- the
\emph{compensation price} $\dt$ throughout this paper --- is the 90th percentile
of South Australian spot prices over the trailing 365 days. Three design features
matter for everything that follows.

First, the payment is effectively \emph{gross}: compensation tracks the unit's
entire directed output valued at $\dt$, not the increment over what the unit
would have produced anyway. On the 271 direction episodes for which
unit-level dollars are observable, compensation fits
$0.95 \times \text{directed MWh} \times \dt$ with $R^{2}=0.99$; a
wedge-over-counterfactual model fits at $R^{2}=0.36$ and its contribution is
driven out entirely once the gross term enters.\footnote{Unit-episode dollars
exist only from October 2023, when AEMO's reporting format changed. The
gross-payment characterisation extends to 2022--2023 by the unchanged
compensation methodology, not by direct measurement.} Because a directed unit's
counterfactual market revenue in these hours is near zero --- directions fire
when prices are low --- the gross structure means a direction pays the unit
roughly $\dt$ per MWh for energy the market values at almost nothing.

Second, the reference price looks backward. A trailing-365-day percentile
embeds last year's scarcity in this year's payment. The 2022 east-coast gas
crisis pushed South Australian spot prices, and therefore $\dt$, to
\$348--350/MWh by mid-2022; $\dt$ then remained at that level until mid-2023 --- while gas prices fell
from \$29.85/GJ (mid-2022) to \$12.59/GJ (early 2023) --- because the crisis
stayed inside the trailing window. The margin
of $\dt$ over the focal units' fuel cost swung from $-\$185$/MWh in April 2022
to $+\$211$/MWh in January 2023 (Figure~\ref{fig:rent}). The scheme spent 2023
paying 2022 scarcity prices into a cheap-fuel world.

Third, the formula payment is a floor, not a ceiling. Clause 3.15.7B entitles a
directed participant to additional compensation --- loss of revenue and net
direct costs above the formula amount, assessed by an independent expert --- so
a direction cannot pay less than the unit's cost of complying. The effective
payment is approximately $\max(\dt, c)$ for unit cost $c$, and the net prize a
direction carries is at most $\max(\dt - c, 0)$.\footnote{Additional-compensation
claims are observable only from October 2023; they appear in 190 of 271
unit-episodes but are immaterial in magnitude, because $\dt$ far exceeded cost
throughout that window and the floor never bound. In the months where the floor
binds --- fuel cost above $\dt$ from April to June 2022 for the Torrens units ---
the channel is unmeasured, and the make-whole characterisation rests on the
rule, not on observed payments. Section~\ref{sec:robustness} reprices the RQ2
dose under this assumption.}

\subsection{The direction regime in South Australia, 2021--2025}

Directions were once near-continuous: by this paper's event record ---
which reproduces AEMO's published share-of-time-directed statistics
(Section~\ref{sec:data}) --- South Australia spent 86\% of the final quarter
of 2021 under an active direction. Four high-inertia synchronous
condensers commissioned by the transmission operator entered service across
late 2021, and the direction regime changed shape: the average megawatts online
under direction fell from 117--196~MW in 2021 to a stable 49--77~MW from 2022
onward, as directions shifted from bulk system-strength procurement to
targeting the marginal synchronous unit needed to complete a minimum
combination. The 2022--2024 sample window used throughout this paper covers
this mature regime: directions are frequent (1,638 direction events), targeted,
and dominated by a small set of units.

Three of those units anchor the analysis: TORRB2, TORRB3, and TORRB4, the
200-MW gas steam units of AGL's Torrens Island B station, the workhorses of the
direction regime and its compensation. Pelican Point (PPCCGT, 478~MW), the
state's most efficient gas plant, serves as the primary comparison unit; it is
directed rarely and runs commercially. Osborne (OSB-AG, 180~MW) is
near-must-run and enters descriptively only.\footnote{Barker Inlet is excluded:
it offers no cheap tranche even in fully competitive intervals, so the
withholding precondition fails by construction.}

The regime has since ended. In September 2025 the minimum synchronous
requirement was cut to one unit, with a target of zero once the second stage of
the Project EnergyConnect interconnector completes. The 2022--2024 window is
therefore not an arbitrary sample: it brackets the years in which a
backward-looking gross-payment scheme, a binding synchronous requirement, and a
fleet of otherwise-unprofitable steam units coexisted. Section~\ref{sec:discussion}
returns to the exit.
